\documentclass[../../../../main.tex]{subfiles}

\begin{document}

\section*{第1問}

\begin{proof}[解]
\textbf{(1)} $O$を原点とし，$\vec a=\overrightarrow{OA}$，$\vec b=\overrightarrow{OB}$，$\vec c=\overrightarrow{OC}$（$|\vec a|=|\vec b|=|\vec c|=1$，$\vec a\cdot\vec b=\vec b\cdot\vec c=\vec c\cdot\vec a=\dfrac12$）とおく。$D$は$AB$の中点，$E$は$OC$の中点だから
\begin{align*}
\overrightarrow{AC}=\vec c-\vec a,\qquad\overrightarrow{DE}=\frac{\vec c}2-\frac{\vec a+\vec b}2=\frac12(\vec c-\vec a-\vec b).
\end{align*}
よって
\begin{align*}
\overrightarrow{AC}\cdot\overrightarrow{DE}&=\frac12\bigl\{(\vec c-\vec a)\cdot(\vec c-\vec a)-(\vec c-\vec a)\cdot\vec b\bigr\}=\frac12\bigl\{|\vec c-\vec a|^2-\vec b\cdot\vec c+\vec a\cdot\vec b\bigr\}.
\end{align*}
$|\vec c-\vec a|^2=|\vec c|^2-2\vec a\cdot\vec c+|\vec a|^2=1-1+1=1$，$\vec b\cdot\vec c=\vec a\cdot\vec b=\dfrac12$だから
\begin{align*}
\overrightarrow{AC}\cdot\overrightarrow{DE}=\frac12\Bigl(1-\frac12+\frac12\Bigr)=\frac12.
\end{align*}

\textbf{(2)} 目の積が10の倍数（$2$の倍数かつ$5$の倍数）にならない事象を考える。$B$:「3回とも偶数の目が出ない（目が$1,3,5$のみ）」，$C$:「3回とも$5$が出ない（目が$1,2,3,4,6$のみ）」，$A=B\cap C$:「3回とも目が$1,3$のみ」とすると，求める事象の余事象は$B\cup C$である。
\begin{align*}
P(A)=\Bigl(\frac26\Bigr)^3=\Bigl(\frac13\Bigr)^3,\qquad P(B)=\Bigl(\frac36\Bigr)^3=\Bigl(\frac12\Bigr)^3,\qquad P(C)=\Bigl(\frac56\Bigr)^3.
\end{align*}
包除原理より$P(B\cup C)=P(B)+P(C)-P(B\cap C)=P(B)+P(C)-P(A)$だから，求める確率は
\begin{align*}
P(\overline{B\cup C})=1-P(B)-P(C)+P(A)=1-\Bigl(\frac12\Bigr)^3-\Bigl(\frac56\Bigr)^3+\Bigl(\frac13\Bigr)^3=1-\frac18-\frac{125}{216}+\frac1{27}=\frac{216-27-125+8}{216}=\frac{72}{216}=\frac13.
\end{align*}

（別解：直接数え上げる。3個の目の中に「5の倍数」（=5）と「2の倍数」（=2,4,6）が少なくとも1つずつ含まれる場合を数える。5の倍数の個数と2の倍数の個数の内訳で場合分けすると，$3!\cdot1\cdot3\cdot2=36$，${}_3C_1\cdot3\cdot1=9$，${}_3C_1\cdot1\cdot3=27$の3通りの内訳があり，合計$36+9+27=72$通り。よって確率は$72/6^3=72/216=1/3$。）
\end{proof}

\newpage

\section*{第2問}

\begin{proof}[解]
\textbf{(1)} $A=\displaystyle\sum_{n=0}^{99}3^n=\frac{3^{100}-1}{3-1}=\frac12(3^{100}-1)=\frac{3^{100}}2-\frac12$とおく。$3^{100}/2$は整数ではないが，桁数を考える上で$A$は$B:=3^{100}/2$とほぼ同じ大きさであり，$3^{100}/2$が$10^\ell$（$\ell\in\mathbb Z$）の形になることはないから，$A$の桁数は$B$の桁数と一致する。$B$が$m$桁とすると$10^{m-1}\le B<10^m$。常用対数をとって
\begin{align*}
m-1\le\log_{10}B<m,\qquad\log_{10}B=100\log_{10}3-\log_{10}2=47.71-\log_{10}2.\tag{①}
\end{align*}
$\log_{10}x$は単調増加関数であり$0=\log_{10}1<\log_{10}2<\log_{10}3=0.4771$だから
\begin{align*}
47.71-0.4771<47.71-\log_{10}2<47.71-0,\quad\text{すなわち}\quad47.2329<\log_{10}B<47.71
\end{align*}
となり，これは常に$47<\log_{10}B<48$の範囲にあるから，①をみたす$m$は$m=48$のみ。よって$A$は$\boxed{48}$桁である。

（別解：$0.4771<0.71<0.4771\times2=\log_{10}9$の各辺に$47$を加えて$\log_{10}10^{47}+\log_{10}3<47.71<\log_{10}10^{47}+\log_{10}9$，すなわち$3\cdot10^{47}<3^{100}<9\cdot10^{47}$（$\log_{10}x$は単調増加）。よって$\dfrac32\cdot10^{47}-\dfrac12<A<\dfrac92\cdot10^{47}-\dfrac12$となり，これより$A$は$48$桁。）

\textbf{(2)} $m^2\le n<(m+1)^2$（$m\in\mathbb N$）のとき$m\le\sqrt n<m+1$より$[\sqrt n]=m$。したがって，この範囲で$n$が$m$で割り切れるのは
\begin{align*}
m=1\text{のとき}&:\ n=1,2,3\ \text{の}3\text{通り},\\
m\ge2\text{のとき}&:\ n=m^2,\ m(m+1),\ m(m+2)\ \text{の}3\text{通り}
\end{align*}
である（区間$[m^2,m^2+2m]$内の$m$の倍数はこの3つに限る）。$10000=100^2$だから，$n=10000$（$m=100$のとき$n=m^2$）も条件をみたす。$m=1,\dots,99$については各々$3$通りずつ，$m=100$については$n=10000$の$1$通りのみ（範囲が$n\le10000$で打ち切られるため）。よって求める個数は
\begin{align*}
\sum_{m=1}^{99}3+1=297+1=298.
\end{align*}
\end{proof}

\newpage

\section*{第3問}

\begin{proof}[解]
$f(x)=x^3-3x^2+2x$とおく。$f'(x)=3x^2-6x+2$。

\textbf{(1)} $x\neq0$のとき，$f(x)=ax\iff x^2-3x+2=a$（①）。①が実数解を持てばよく，左辺$=(x-\frac32)^2-\frac14$の値域は$\ge-\dfrac14$なので，①が解を持つ条件は
\begin{align*}
a\ge-\frac14.
\end{align*}

\textbf{(2)} $2\le a$のとき$S(a)$は単調に増加するので，$-\dfrac14\le a\le2$の範囲で考えれば十分である。①の2実解を$\alpha\le\beta$とすると，解と係数の関係より$\alpha+\beta=3$，$\alpha\beta=2-a$。$a<2$では$\alpha\beta>0$かつ$\alpha+\beta=3>0$より$0<\alpha\le\beta$（$a=2$では$\alpha=0$）。
\begin{align*}
\alpha=\frac{3-\sqrt{4a+1}}2,\qquad\beta=\frac{3+\sqrt{4a+1}}2.
\end{align*}
$g(x)=f(x)-ax=x(x-\alpha)(x-\beta)$は，$0<x<\alpha$で正，$\alpha<x<\beta$で負だから，$C$と$l$で囲まれる面積は
\begin{align*}
S(a)=\int_0^\alpha g(x)dx-\int_\alpha^\beta g(x)dx=\int_0^\alpha f(x)dx-\int_0^\alpha ax\,dx-\int_\alpha^\beta f(x)dx+\int_\alpha^\beta ax\,dx.\tag{①}
\end{align*}
$\alpha,\beta$は$a$の関数であることに注意して両辺を$a$で微分する。$f(\alpha)=a\alpha$，$f(\beta)=a\beta$（①の定義より）を用いると
\begin{align*}
\frac d{da}\int_0^\alpha f(x)dx&=f(\alpha)\alpha'=a\alpha\alpha',\\
\frac d{da}\int_\alpha^\beta f(x)dx&=f(\beta)\beta'-f(\alpha)\alpha'=a(\beta\beta'-\alpha\alpha'),\\
\frac d{da}\int_0^\alpha ax\,dx&=\int_0^\alpha x\,dx+a\alpha\alpha'=\frac12\alpha^2+a\alpha\alpha',\\
\frac d{da}\int_\alpha^\beta ax\,dx&=\int_\alpha^\beta x\,dx+a(\beta\beta'-\alpha\alpha').
\end{align*}
よって①の両辺を$a$で微分して
\begin{align*}
S'(a)&=a\alpha\alpha'-\Bigl(\frac12\alpha^2+a\alpha\alpha'\Bigr)-a(\beta\beta'-\alpha\alpha')+\int_\alpha^\beta x\,dx+a(\beta\beta'-\alpha\alpha')\\
&=-\frac12\alpha^2+\int_\alpha^\beta x\,dx=-\frac12\alpha^2+\frac12(\beta^2-\alpha^2)=\frac12(\beta^2-2\alpha^2)=\frac12(\beta-\sqrt2\alpha)(\beta+\sqrt2\alpha).\tag{②}
\end{align*}
$0\le\alpha\le\beta$より$\beta+\sqrt2\alpha\ge0$（等号は$\alpha=\beta=0$のときのみ）だから，$S'(a)$の符号は$\beta-\sqrt2\alpha$の符号で決まる。$A=\sqrt{4a+1}$（$0\le A\le3$）とおくと$\alpha=\frac{3-A}2$，$\beta=\frac{3+A}2$だから
\begin{align*}
\beta-\sqrt2\alpha=\frac12\bigl\{(3+A)-\sqrt2(3-A)\bigr\}=\frac{1+\sqrt2}2\bigl\{A-3(3-2\sqrt2)\bigr\}
\end{align*}
（$3(1-\sqrt2)/(-(1+\sqrt2))$の整理により$3(\sqrt2-1)/(1+\sqrt2)=3(\sqrt2-1)^2=3(3-2\sqrt2)$を用いた）。よって$S'(a)$は$A=3(3-2\sqrt2)=9-6\sqrt2$を境に符号が$-$から$+$に変わる。このとき
\begin{align*}
4a+1=A^2=(9-6\sqrt2)^2=81-108\sqrt2+72=153-108\sqrt2
\end{align*}
より
\begin{align*}
a=\frac{152-108\sqrt2}4=38-27\sqrt2.
\end{align*}
（$38-27\sqrt2\approx-0.18$は$[-1/4,2]$の範囲内。）$S'(a)$は$-\dfrac14\le a<38-27\sqrt2$で負，$38-27\sqrt2<a\le2$で正だから，$S(a)$は
\begin{align*}
a=38-27\sqrt2
\end{align*}
で最小となる。
\end{proof}

\newpage

\section*{第4問}

\begin{proof}[解]
\begin{align*}
a_1=\frac1{n(n+1)},\qquad a_{k+1}=-\frac1{k+n+1}+\frac nk\sum_{i=1}^ka_i\qquad(k=1,2,\dots).\tag{①}
\end{align*}

\textbf{(1)} ①より
\begin{align*}
a_2&=-\frac1{n+2}+n\cdot\frac1{n(n+1)}=-\frac1{n+2}+\frac1{n+1}=\frac1{(n+1)(n+2)},\\
a_3&=-\frac1{n+3}+\frac n2\Bigl(\frac1{n(n+1)}+\frac1{(n+1)(n+2)}\Bigr)=-\frac1{n+3}+\frac n2\cdot\frac1{n+1}\cdot\frac{2(n+1)}{n(n+2)}=-\frac1{n+3}+\frac1{n+2}=\frac1{(n+2)(n+3)}.
\end{align*}

\textbf{(2)} $a_k=\dfrac1{(n+k-1)(n+k)}$であることを数学的帰納法で示す。$k=1$のときは$a_1=\dfrac1{n(n+1)}$と一致し成立。$k\le m$（$m\in\mathbb N$）なる全ての$k$で成立すると仮定する。$S_m=\displaystyle\sum_{i=1}^ma_i$は望遠鏡和で
\begin{align*}
S_m=\sum_{i=1}^m\Bigl(\frac1{n+i-1}-\frac1{n+i}\Bigr)=\frac1n-\frac1{n+m}.
\end{align*}
①より
\begin{align*}
a_{m+1}=-\frac1{n+m+1}+\frac nm\Bigl(\frac1n-\frac1{n+m}\Bigr)=-\frac1{n+m+1}+\frac nm\cdot\frac m{n(n+m)}=-\frac1{n+m+1}+\frac1{n+m}=\frac1{(n+m)(n+m+1)}
\end{align*}
となり，これは$k=m+1$でも公式が成立することを示す。よって数学的帰納法により，すべての$k$で
\begin{align*}
a_k=\frac1{(n+k-1)(n+k)}.
\end{align*}

\textbf{(3)} $b_n=\displaystyle\sum_{k=1}^n\sqrt{a_k}$とする。
\begin{align*}
\sum_{k=1}^n\frac1{n+k}=\frac1n\sum_{k=1}^n\frac1{1+\frac kn}\xrightarrow[n\to\infty]{}\int_0^1\frac{dx}{1+x}=\log2,\qquad
\sum_{k=1}^n\frac1{n+k-1}=\frac1n\sum_{k=1}^n\frac1{1+\frac{k-1}n}\xrightarrow[n\to\infty]{}\int_0^1\frac{dx}{1+x}=\log2
\end{align*}
（区分求積法）。一方，$(n+k-1)^2<(n+k-1)(n+k)<(n+k)^2$（$n+k-1<n+k$より）から逆数をとって
\begin{align*}
\frac1{(n+k)^2}<a_k<\frac1{(n+k-1)^2}\quad\therefore\ \frac1{n+k}<\sqrt{a_k}<\frac1{n+k-1}
\end{align*}
となり，$k=1,\dots,n$について和をとると
\begin{align*}
\sum_{k=1}^n\frac1{n+k}<b_n<\sum_{k=1}^n\frac1{n+k-1}.
\end{align*}
両端はともに$n\to\infty$で$\log2$に収束するから，はさみうちの原理より
\begin{align*}
\lim_{n\to\infty}b_n=\log2.
\end{align*}
\end{proof}

\newpage

\section*{第5問}

% (原稿において第5問の解答は空欄)
\newpage

\end{document}
